\documentclass[12pt]{article}

\usepackage[dvips,letterpaper,margin=0.75in,bottom=0.5in]{geometry}
\usepackage{cite}
\usepackage{slashed}
\usepackage{graphicx}
\usepackage{amsmath}
\usepackage{amssymb}
\usepackage{braket}
\usepackage{pythonhighlight}
\usepackage{mathtools}
\begin{document}

\newcommand{\ihbar}{\ensuremath{i \hbar}}
\newcommand{\Pss}{\ensuremath{\Psi^*}}
\newcommand{\dPsidt}{\ensuremath{ \frac{\partial \Psi}{\partial t} }}
\newcommand{\dPsidx}{\ensuremath{ \frac{\partial \Psi}{\partial x} }}
\newcommand{\ddPsidx}{\ensuremath{ \frac{\partial^2 \Psi}{\partial x^2} }}
\newcommand{\dPssdt}{\ensuremath{ \frac{\partial \Psi^*}{\partial t} }}
\newcommand{\dPssdx}{\ensuremath{ \frac{\partial \Psi^*}{\partial x} }}
\newcommand{\ddPssdx}{\ensuremath{ \frac{\partial^2 \Psi^*}{\partial x^2} }}

\newcommand{\dphidt}{\ensuremath{ \frac{d \phi}{dt} }}
\newcommand{\dpsidx}{\ensuremath{ \frac{d \psi}{dx} }}
\newcommand{\ddpsidx}{\ensuremath{ \frac{d^2 \psi}{dx^2} }}


\title{PHY 115L \\ Variational Methods \\ using Fourier Series }
\author{Michael Mulhearn}

\maketitle



\section{Notation}

We'll reserve $k$ for the wave number of a particle.  We will use $m$
and $n$ as dummy indices.  So we will use $M$ (capital) for the mass
of a particle to avoid confusion with the index $m$.  We'll use $N$
for the number of samples.  We'll use $a$ for the period of a
function.  Note that the numpy FFT library uses different notation,
taking $k$ as the index for Fourier coefficients (our $m$) and $n$ for
the number of samples (our $N$).  The HW1 write-up has been updated to
be consistent with this notation in case you want to refer back to
anything there.  In our notation, the wave number associated with the
Fourier index $m$ is given by:
\begin{equation}
k = \frac{2 \pi m}{a}
\end{equation}

\section{TISE and System of Units}

\noindent
We will use a variational approach to find the ground state
$\Psi_0(x)$ of the Time-Independent Schr\"odinger Equation (TISE):
$$\hat{H} \; \psi_0(x) = E_0 \; \psi_0(x)$$
where the Hamiltonian operator is defined by:
\begin{equation}
\hat{H} = -\frac{\hbar^2}{2 M} \; \frac{d^2}{dx^2} \, + \, V(x)
\end{equation}
We will consider potentials with a characteristic length scale $L$.
We will be considering {\bf non-periodic} wave function
(e.g. solutions to the harmonic oscillator).  To use the Fourier
series approach in this case, we take the period $a$ to be large
compared to $L$, and use the fact that the wave function approaches
zero to approximate a periodic function.

Based on our characteristic length scale $L$, a characteristic energy is:
\begin{equation}
\epsilon \equiv \frac{\hbar^2}{2 M L^2}
\end{equation}
Note well:  your program should not be calculating a numerical value for
$\epsilon$: that is the unit of energy for your program!  For example
$V(x) = 2$ in your program represents $V(x) = 2 \epsilon$
physically.  Using a system of units like this ensures that the
numbers we are calculating will be near $1$, which minimizes round-off
error.

With these units, we have:
\begin{equation}
\frac{\hat{H}}{\epsilon} = -L^2\frac{d^2}{dx^2} + \frac{V(x)}{\epsilon}
\end{equation}

\section{The Variational Method of Finding the Ground State}

Our technique is based on the observation that:
$$\braket{\psi|\hat{H}|\psi} \geq E_0$$
for all $\psi$.  Which you can see by expanding $\ket{\psi}$ in the complete set of energy eigenstates:
$$\ket{\psi} = \sum c_m \ket{E_m}$$
so:
$$\braket{\psi|\hat{H}|\psi} = \sum |c_m|^2 E_m = |c_0|^2 E_0 + \sum_{m>0} |c_m|^2 E_m$$
but:
$$\sum |c_m|^2 = 1 \;\; \implies \;\; |c_0|^2 = 1 - \sum_{m>0} |c_m|^2$$
so:
$$\braket{\psi|\hat{H}|\psi} = E_0 + \sum_{m>0} |c_m|^2 (E_m - E_0)$$
but as $E_m \geq E_0$ we have:
$$\braket{\psi|\hat{H}|\psi} \geq E_0$$

The variational technique starts with a first guess for the wave function $\psi$, for which we calculate:
$$E=\braket{\psi|\hat{H}|\psi}$$
During each iteration, we vary the wave function randomly to
obtain $\psi'$, and calculate:
$$E'=\braket{\psi'|\hat{H}|\psi'}$$
If $E' < E$, then we update $\psi$ to $\psi'$ and $E$ to $E'$.
Otherwise, we leave $\psi$ and $E$ as is.  Then we move on to the next
iteration.

\section{Fourier Series}

One can apply the variational principle directly to the wave function.
We approximate the continuous wave function as a set of y-values for
each of $N$ x-values.  To estimate the kinetic energy we take a
numerical second derivative.  However, to approximate a continuous
function and its second derivative, we need to use lots of $x$-values
(order 100).  And we have to vary each point, randomly, one by one.
And so the convergence of the variational method with this approach is
very slow, even for a relatively simple problem.

We will instead use a Fourier series approach, taking advantage of the
FFT library in numpy.  The Fourier series in many cases converges
remarkably fast.  We will only need $N=16$ to handle the harmonic
oscillator potential with high accuracy, for example.  And calculating
the expectation value of the kinetic energy is very easy in frequency
space, much easier than numerically calculating a second derivative.

We will be considering a Fourier series expansion of the wave function:
\begin{equation}
\label{eqn:fs}
\psi(x) = \sqrt{\frac{1}{a}} \sum_{m=-\infty}^{+\infty} c_m \exp(i\frac{2\pi m x}{a})
\end{equation}
where our normalization here is such that:
$$\sum_{m=-\infty}^{+\infty} |c_m|^2 = 1$$
As we will be considering stationary states of a 1-D potential, we can take $\psi(x)$ to be real, and so:
$$c_{-k} = c_k^*$$
During each iteration we will calculate the expectation value:
$$\braket{\psi|\hat{H}|\psi} = \braket{\psi|K|\psi} + \braket{\psi|V(x)|\psi}$$
where $K$ refers to the kinetic energy.

The Fourier series is an expansion in planes waves with momentum $p = \hbar k$ and wavenumbers:
$$k = \frac{2 \pi m}{a}$$
And so the kinetic energy $K(m)$ of a particle described by the plane wave for coefficient $c_m$ is:
$$ K(m) = \frac{p^2}{2M} = \frac{(\hbar k)^2}{2M}
= \frac{(2 \pi \hbar m)^2}{2Ma^2} = \left( \frac{\hbar^2}{2ML^2} \right) (2 \pi m)^2 \left( \frac{L}{a}\right)^2$$
And so:
$$\frac{K(m)}{\epsilon} = (2 \pi m)^2 \left( \frac{L}{a}\right)^2 $$
From the Fourier series coefficients we can therefore easily calculate the expectation value of the kinetic energy as an average over the kinetic energy of each plane wave:
$$\braket{\psi|K|\psi} = \epsilon \sum_m (2\pi m)^2 |c_m|^2 \left( \frac{L}{a}\right)^2 $$
The inverse discrete Fourier series, acting on these $n$ coefficients, will provide us with $\psi(x_m)$ for the $n$ $x$-values located at:
$$x_m = x_0 + m \Delta a$$
with
$$\Delta a \coloneqq \frac{a}{n}$$
We will use these to calculate:
$$\braket{\psi|V(x)|\psi} = \sum_m V(x_m) \, |\psi(x_m)|^2 \, \Delta a $$
Keep in mind what is changing and what is constant.  We'll be updating Forier coefficients ($c_m$) and that will change the values of $\psi(x_m)$ returned by the IFFT.  But the IFFT really only sees these as $\psi_m$.  The $x_m$ are fixed, and so your $V(x_m)$ are also fixed.

\section{Real FFT in Numpy}

We can always take our stationary solutions to 1-D potentials as real-valued.  This implies that:
$$c_{-m} = c_{m}^*$$
which implies that $c_0$ is real valued.  As we will be randomly
varying these coefficients, we want to take full advantage of these
constraints.  Therefore, we will be using the {\bf real} fast Fourier
transform (RFFT), which has some different behavior from the fast
fourier transform for a complex valued function that we need to sort
out.

Given $N$ data points, the RFFT provides coefficients $c_m$ for
$$m = 0,1,2,\cdots,m_{\rm MAX}$$
where:
$$m_{\rm MAX} = (N//2) + 1$$
and the $//$ refers to integer divide, e.g.:
$$5 // 2 = 2$$
Notice that the same number of coefficients are provided for
e.g. $n=4$ and $n=5$. {\bf Important:} because of this ambiguity, you
must specify the number of original data points via the option {\tt n=N}
when calling the IFFT routine.  Otherwise, the number of sample points
and there implicit locations may not be consistent with what you
started with!

The fast Fourier transform, as implemented in numpy, when using the {\tt norm=``ortho''} option, is:
$$A_m = \sqrt{\frac{1}{n}} \sum_{j=0}^{n-1} a_j \exp \left( -i \frac{2 \pi m j}{n}\right)$$
where the $a_j$ are the sampled function values:
$$a_m = \psi(x_m) \psi(x_0 + \Delta a m)$$
with inverse:
$$a_m = \sqrt{\frac{1}{n}} \sum_{j=0}^{n-1} A_j \exp \left( i \frac{2 \pi m j}{n}\right)$$
This differs from our (physics) based normalization which puts:
$$\sum_{m=-\infty}^{+\infty} |c_m|^2 = 1$$
In the previous assignment, we already worked out the conversion between our
physics-based coefficients ($c_m$), and those returned by the tool
($A_m$), as:
\begin{equation}
\label{eqn:am}
A_m \equiv \frac{c_m}{\sqrt{a/N}} \exp\left(i \frac{2\pi m x_0}{a}\right)
\end{equation}
Keep in mind that while the tool provides $A_m$ for $m=0,1,2,\ldots,m_{\rm MAX}$, it is with the understanding that one obtains the coefficiets for $m=-1,-2,\ldots,-m_{\rm MAX}$ via:
$$A_{-m} = A_{m}^*$$
So we have:
\begin{eqnarray*}
  \sum_{m=\infty}^{m=+\infty} |c_m|^2
  &=& \left( \frac{a}{N} \right) \sum_{m=-m_{\rm MAX}}^{m_{\rm MAX}} |A_m|^2 \\
  &=& ( \Delta a ) \left( |A_0|^2 + 2 \sum_{m=0}^{m_{\rm MAX}} |A_m|^2 \right) \\
\end{eqnarray*}
So we can state the normalization of the Fourier coefficients $A_k$ as:
\begin{equation}
  \left(|A_0|^2 + 2 \cdot \sum_{m=1}^{m_{\rm MAX}} |A_m|^2 \right) \Delta a = 1
\end{equation}
Be careful about that factor of two.  It doesn't apply to the $m=0$
term, and so you can't simply sum across the all $|A_m|^2$!

To see how to reconstruct the continuous real function $\psi(x)$ from
the Fourier coefficients $A_m$, plug Equation~\ref{eqn:am} into
Equation~\ref{eqn:fs} to obtain (with some work):
\begin{equation}
  \psi(x) = \frac{1}{\sqrt{N}} \left( A_0 + \sum_{m=1}^{m_{\rm MAX}}
  \left\{
2 \, {\rm Re}(A_m) \, \cos\left( \frac{2 \pi m (x-x_0)}{a}\right)
- 2 \, {\rm Im}(A_m) \, \sin\left( \frac{2 \pi m (x-x_0)}{a}\right) \right\} \right)
\end{equation}

\section{Numerical Algorithm}

Suppose we are trying to find the ground state of a potential and it
is {\bf even}.  This means that all of the Fourier coefficients are real (make sure you understand this!).
Our algorithm starts with some preliminaries:
\begin{itemize}
\item Initialize a numpy array of your best-so-far Fourier
  coefficients ($A_m$) by filling them with random numbers, but be sure that
  each coefficient is real (zero imaginary component) as appropriate
  for a real function.
\item Normalize your best-so-far Fourier coefficients appropriately.
\item Calculate the expectation value for the total energy $E$ of the best-so-far coefficients.
\end{itemize}
Then we enter an iterative process, which we run some number of times.  At each iteration:
\begin{itemize}
 \item Make a copy of best-so-far Fourier coefficients (call it $A'_m$).
 \item Choose a random $m$ value and change $A'_m$ by a random amount uniform in $\pm \Delta A$.
 \item Normalize $A'_m$ appropriately.
 \item Calculate total energy $E'$ of the coefficients $A'_m$.
 \item If $E' < E$, then update $A_m$ and $E$ to $A'_m$ and $E$.  Otherwise, abandon $A'_m$ and enter the next iteration.
\end{itemize}

Make sure you understand how numpy arrays work with respect to
copying.  For example, it is very easy to introduce a bug that changes
your best-so-far coefficients ($A_m$) when you only intended to change
your trial coefficients ($A'_m$).  Write a simple little test code to
check your approach to copying, instead of trying to test your
understanding deep in the interior of a complicated algorithm.

For precise results, it is helpful to use a process called simulated
annealing, where we initially use a large value of $\Delta A$, so that
the wave function changes rapidly, then reduce the value of $\Delta A$
for fine tuning as we get closer to the ground state.

\section{The Simple Harmonic Oscillator}

In our system of units, the simple harmonic oscillator has potential:
$$\frac{V(x)}{\epsilon} = \left(\frac{x}{L}\right)^2$$
Remember, that in your program, this would be programmed as:
\begin{python}
def V(x):
  return x**2
\end{python}
because we set $L=1$ and $\epsilon=1$ and leave them out of the
program entirely.  The ground state is:
$$\psi_0 \left( \frac{x}{L} \right) = \frac{1}{\pi^\frac{1}{4}}
\exp\left( -\frac{x^2}{2 L^2} \right)$$
and the first excited state:
$$\psi_1 \left( \frac{x}{L} \right) = \frac{\sqrt{2}}{\pi^\frac{1}{4}} x \exp\left( -\frac{x^2}{2 L^2} \right) $$
And the energy of each state is:
$$\frac{E_n}{\epsilon} = 2 n + 1$$
Last of all, we know that for this potential:
$$\braket{V} = \braket{K}$$
Analytic solutions like the Harmonic oscillator are absolute gems for
developing numerical algorithms.  It is a much much more difficult problem
to debug an algorithm when you do not know the correct answer!

\section{Homework Problems}

\noindent
{\bf Please submit one PDF file (including output) and one python
  notebook file (.ipynb) of your completed assignment.}  If using
C/C++, submit text or PDF of your output and the source code.

\noindent
{\bf Problem 1:} For the harmonic oscillator, and using their known
analytic functions, plot $V(x)$, $\psi_0(x)$, and $\psi_1(x)$ versus
$x$ in our system of units. For better visiblity, offset the
$y$-values by the energy of each state.  So e.g. plot $\psi_0(x)=0$ at
$V(x)/\epsilon=1$.  See Griffith's Figure 2.7(a) which you are
reproducing.\\[3pt]

\noindent
{\bf Problem 2:} Demonstrate mastery of the real fast fourier
transform (RFFT) and its inverse (IRFFT) by doing the following.
\begin{itemize}
 \item Take $N=16$, $a=10$, and sample $\psi_0(x)$ appropriately. 
 \item Plot A: Plot your sample points (as data points) on top of the
   continuous function (fine binned line) in the same plot.
 \item Plot B:  Use the RFFT to obtain the Fourier coefficients $A_m$ and plot them versus their index $m$.
 \item Plot C:  Use the IRFFT to reproduce the sample points and plot them as in Plot A. If you have done this all correctly, then Plot C should reproduce Plot A. \\[3pt]
\end{itemize}

\noindent
{\bf Problem 3:} Repeat Problem 2 for $\psi_1(x)$.\\[3pt]

\noindent
{\bf Problem 4:} Write a utility function to check the normalization of the $A_k$ and test it on the Fourier coefficients obtained for $\psi_0(x)$ and $\psi_1(x)$.\\[3pt]

\noindent
{\bf Problem 5:} Write a utility function to measure the expectation value for the kinetic energy from the Fourier coefficients $A_k$.   Validate it on $\psi_0(x)$ and $\psi_1(x)$.\\[3pt]

\noindent
{\bf Problem 6:} Write a utility function to measure the expectation
value for the potential energy from the Fourier coefficients $A_k$ and
the sampled potential $V(x_m)$.  This will use the IRFFT.  Validate it
on $\psi_0(x)$ and $\psi_1(x)$.\\[3pt]

\noindent
{\bf Problem 7:} Use the variational method to calculate the ground state of the harmonic oscillaor.  Report the ground state energy and plot the wave function.\\[3pt]

\noindent
{\bf Problem 8:} Use the variational method to calculate the first
excited state of the harmonic oscillator by modifying the algorithm
appropriately.  (Hint: the first excited state is an {\bf odd} wave
function.)  Report the ground state energy and plot the wave
function.\\[3pt]

\noindent
{\bf Problem 9:} For problem 7, plot your starting guess, your final wave function, and several steps along the way, all in the same plot.\\[3pt]

\end{document}
